% Generated deterministically from the audited Markdown source.
% Responsible author: Carptopus.
\documentclass[11pt]{article}
\usepackage[a4paper,margin=27mm]{geometry}
\usepackage[T1]{fontenc}
\usepackage[utf8]{inputenc}
\usepackage{lmodern}
\usepackage{microtype}
\usepackage{amsmath,amssymb,amsthm,mathtools}
\usepackage{xcolor}
\usepackage{fvextra}
\usepackage[colorlinks=true,linkcolor=blue!55!black,citecolor=blue!55!black,urlcolor=blue!65!black]{hyperref}
\usepackage{xurl}
\usepackage{fancyhdr}

\newtheorem{theorem}{Theorem}
\newtheorem{lemma}[theorem]{Lemma}
\numberwithin{equation}{section}
\pagestyle{fancy}
\fancyhf{}
\fancyhead[L]{Carptopus}
\fancyhead[R]{Weak cube saturation}
\fancyfoot[C]{\thepage}
\setlength{\headheight}{14pt}
\setlength{\parindent}{0pt}
\setlength{\parskip}{0.42em}
\setlength{\emergencystretch}{4em}
\urlstyle{same}

\title{Small exact values and a dense-anchor bound\\for weak cube saturation}
\author{Carptopus\\\href{mailto:carptopus@163.com}{carptopus@163.com}}
\date{Manuscript, 2 September 2026}

\hypersetup{
  pdftitle={Small exact values and a dense-anchor bound for weak cube saturation},
  pdfauthor={Carptopus},
  pdfsubject={Exact weak saturation values for Q3 and a dense-anchor recurrence},
  pdfkeywords={weak saturation, hypercube, Q3, complete graph, exhaustive enumeration, dense-anchor gluing}
}

\begin{document}
\maketitle

\begin{center}
\fcolorbox{black!35}{black!4}{\parbox{0.88\textwidth}{\centering\small
Internally verified candidate manuscript, version 0.1-beta. The mathematical claims and the
complete manuscript have passed project-internal zero-trust audits; external review is pending.}}
\end{center}

\begin{abstract}


We study the weak saturation number of the three-dimensional cube $Q_3$ in the complete graph. By
normalizing the first activated copy of $Q_3$, we determine

\[
\mathrm{wsat}(K_8,Q_3)=15,\qquad
\mathrm{wsat}(K_9,Q_3)=16,\qquad
\mathrm{wsat}(K_{10},Q_3)=18.
\]

The value at $n=9$ corrects a reported value of 17 and disproves the associated conjecture
$\mathrm{wsat}(K_n,Q_3)=2n-1$. We also give an explicit 28-edge weakly saturated graph on 11
vertices containing a complete seven-vertex anchor. It yields

\[
\mathrm{wsat}(K_{N+4},Q_3)\leq\mathrm{wsat}(K_N,Q_3)+7
\]

and hence $\mathrm{wsat}(K_n,Q_3)\leq\lfloor(7n+2)/4\rfloor$ for $n\geq9$. Combined with a
known combinatorial lower bound, this gives $11/7\leq c_{Q_3}\leq7/4$.

\end{abstract}

\section{Weak saturation and the first cube}


A spanning subgraph $H\subseteq K_n$ is weakly $Q_3$-saturated if the missing edges can be added
one at a time so that each added edge completes a new copy of $Q_3$. Let
$\mathrm{wsat}(K_n,Q_3)$ be the minimum initial number of edges.
We use the standard convention that the initial graph need not be $Q_3$-free. Our explicit upper
witnesses satisfy the stronger free condition, while the lower-bound searches do not assume it;
therefore the values below are unchanged under the stricter convention.

Consider the first edge in any valid activation order. Before it is added, the initial graph already
contains all eleven other edges of a copy of $Q_3$. Since $Q_3$ is edge-transitive, this first
$Q_3-e$ witness may be fixed, up to isomorphism, on eight labelled vertices. Therefore an exact
search at a proposed edge count only needs to enumerate choices outside one fixed 11-edge core.

\begin{theorem}


\[
\mathrm{wsat}(K_8,Q_3)=15,\qquad
\mathrm{wsat}(K_9,Q_3)=16,\qquad
\mathrm{wsat}(K_{10},Q_3)=18.
\tag{1}
\]

\end{theorem}

\begin{proof}


For $n=8$, exhaustive enumeration of the normalized 11-edge core plus zero through three
additional edges rules out edge counts 11 through 14; a 15-edge witness closes to $K_8$.

For $n=9$, all normalized candidates through 15 edges were checked, and none closes. A 16-edge
witness has a 20-step activation certificate closing it to $K_9$.

For $n=10$, any weakly saturated initial graph must be biconnected: adding an edge inside a copy of
$Q_3-e$ cannot merge distinct blocks because $Q_3-e$ is biconnected, while the final graph is
biconnected. Filtering by this necessary condition, all normalized candidates through 17 edges were
exhausted. None closes, while an independently verified 18-edge witness closes to $K_{10}$.

The evidence layers are separate. The normalized candidate enumeration is implemented once by
{\small\path{exact_first_cube_enumeration.py}}; the cube edge orbit and the biconnectivity condition are checked
independently with NetworkX; and the $K_9$, $K_{10}$, $K_{11}/K_7$ and $K_{12}$ upper witnesses
are stored as edge-by-edge certificates and replayed by a second NetworkX implementation. In
particular, the 18-edge $K_{10}$ witness has a 27-step certificate closing it to all 45 edges.
This establishes (1).
\end{proof}


A 2025 PCMI report~\cite{pcmi-q3} gave the correct value 15 at $n=8$ but reported 17 at $n=9$ and conjectured
$2n-1$ for all $n\geq8$. The certified 16-edge graph corrects that value and disproves the
conjecture. We do not identify the unreported program error.

\section{Dense-anchor gluing}


\begin{lemma}


Let $B$ be weakly $Q_3$-saturated on $b$ vertices, and let $A\subseteq V(B)$ have size
$s\geq2$. Put

\[
c(B,A)=|E(B)|-|E(B[A])|.
\]

For every $N\geq\max\{8,s\}$,

\[
\mathrm{wsat}(K_{N+b-s},Q_3)
\leq\mathrm{wsat}(K_N,Q_3)+c(B,A).
\tag{2}
\]

\end{lemma}

\begin{proof}


Identify $A$ with $s$ vertices of a weakly saturated graph on the old $N$ vertices. Add the
$b-s$ new vertices and precisely the edges of $B$ having at least one new endpoint. First close the
old graph to $K_N$. All edges inside $A$ are then available, so the copy of $B$ can be closed to
$K_b$.

Choose distinct anchors $a_1,a_2\in A$. For a new vertex $p$ and an old vertex $x$ still not
adjacent to it, use $p$ as one cube vertex and $x,a_1,a_2$ as its three neighbours. Four further
old vertices complete a cube. Every required edge except $px$ is already present, because the old
vertices form a clique and $pa_1,pa_2$ are present after the $K_b$ closure. Thus all cross edges
may be added. The number of initially paid edges is $c(B,A)$.
\end{proof}


\section{A seven-edge four-vertex extension}


Let $A=\{0,1,\ldots,6\}$ span $K_7$. Add four new vertices $7,8,9,10$, the path

\[
7-8-9-10,
\]

and the four spokes

\[
3-7,\qquad4-8,\qquad0-9,\qquad5-10.
\]

Together with the 21 anchor edges, this is a 28-edge graph $B\subset K_{11}$. An explicit
edge-by-edge certificate verifies that it is weakly $Q_3$-saturated and initially $Q_3$-free.
Here $b=11$, $s=7$, and $c(B,A)=7$. Lemma 2 gives

\[
\mathrm{wsat}(K_{N+4},Q_3)
\leq\mathrm{wsat}(K_N,Q_3)+7
\qquad(N\geq8).
\tag{3}
\]

Using the exact seeds at $n=9,10$, a 19-edge seed at $n=11$, and a certified 21-edge seed at
$n=12$, the four residue classes in (3) yield:

\begin{theorem}


For every $n\geq9$,

\[
\mathrm{wsat}(K_n,Q_3)
\leq\left\lfloor\frac{7n+2}{4}\right\rfloor.
\tag{4}
\]

In particular,

\[
c_{Q_3}:=\lim_{n\to\infty}\frac{\mathrm{wsat}(K_n,Q_3)}n\leq\frac74.
\]

\end{theorem}

Terekhov and Zhukovskii's Theorem~2.1~\cite{tz-combinatorial} uses the costs

\[
(e_1,\ldots,e_7)=(2,4,6,7,9,10,11)
\]

for $Q_3$. Its minimum additive closure has period seven: if $n-8=7q+r$ with
$0\leq r<7$, then

\[
\mathrm{wsat}(K_n,Q_3)\geq11+11q+h_r,
\qquad(h_0,\ldots,h_6)=(0,2,4,6,7,9,10).
\]

The standard gluing/subadditivity argument recorded in their Section 2 gives the existence of
$c_{Q_3}$. The seven-periodic formula follows by checking the seven residue classes in the
min-plus recurrence for the costs $e_i$; taking $q$ parts of size seven and one remainder part
attains the resulting lower bound. Hence the displayed specialization yields
$c_{Q_3}\geq11/7$, and

\[
\frac{11}{7}\leq c_{Q_3}\leq\frac74.
\tag{5}
\]

The lower bound in (5) is prior work, not a contribution of this note.

\section{Reproducibility}


The checks use Python 3.13.11 and NetworkX 3.6.1 from the project environment. From the project
root, the exact lower-bound searches are

\begin{Verbatim}[fontsize=\scriptsize,breaklines=true,breakanywhere=true]
.venv\Scripts\python.exe loops\MATH-DIR-0053\verification\exact_first_cube_enumeration.py --n 8 --maximum-edges 15
.venv\Scripts\python.exe loops\MATH-DIR-0053\verification\exact_first_cube_enumeration.py --n 9 --maximum-edges 16
.venv\Scripts\python.exe loops\MATH-DIR-0053\verification\exact_first_cube_enumeration.py --n 10 --maximum-edges 17 --structural-filter
\end{Verbatim}


Through the last excluded edge count, these enumerate respectively 697, 12,951 and 1,391,842
normalized candidates; in the $n=10$ case the proved biconnectivity filter leaves 301,000 candidates
for closure testing. The complete result files are {\small\path{exact-first-cube-n8.json}},
{\small\path{exact-first-cube-n9.json}}, and {\small\path{exact-first-cube-n10.json}} in
{\small\path{loops/MATH-DIR-0053/verification/results/}}.

The $K_9$, $K_{10}$, $K_{11}/K_7$, and $K_{12}$ upper certificates are independently replayed by

\begin{Verbatim}[fontsize=\scriptsize,breaklines=true,breakanywhere=true]
.venv\Scripts\python.exe loops\MATH-DIR-0053\verification\audit_activation_certificate.py <certificate.json>
\end{Verbatim}


using the four certificate files in the same results directory. The complete command list and the
logical boundary between exhaustive search, structure checks, and certificate replay are recorded
in {\small\path{loops/MATH-DIR-0053/verification/README.md}}.

\section{Scope}


The note does not determine $\mathrm{wsat}(K_n,Q_3)$ for every $n$, nor the exact asymptotic
constant. Overlapping-block gluing is also present in earlier weak-saturation work; the claimed
increment is the concrete complete-anchor lemma, the certified $K_{11}/K_7$ block, the resulting
$7/4$ upper bound, and the corrected exact values.

An attempted extension did not improve either side of (5). In particular, the standard
single-interface reuse of the certified 21-edge $K_{12}$ seed costs two edges per new vertex, and a
uniform sparsity-matroid lower-bound scheme is capped at slope $11/7$. These are stopping results,
not stronger bounds.

\begin{thebibliography}{9}
\footnotesize
\raggedright

\bibitem{tz-combinatorial}
N.~Terekhov and M.~Zhukovskii,
\emph{Weak saturation in graphs: a combinatorial approach},
Journal of Combinatorial Theory, Series B \textbf{172} (2025), 146--167.
\href{https://arxiv.org/abs/2305.11043}{arXiv:2305.11043}.

\bibitem{tz-rank}
N.~Terekhov and M.~Zhukovskii,
\emph{Weak saturation rank: a failure of linear algebraic approach to weak saturation},
\href{https://arxiv.org/abs/2405.17857}{arXiv:2405.17857}.

\bibitem{pcmi-q3}
T.~Trakulthongchai, D.~Xu and K.~Zhu,
\emph{Weak saturation number of the 3-cube in the complete graph},
Park City Mathematics Institute, Summer 2025.
\url{https://www.ias.edu/sites/default/files/Trakulthongchai_Xu_Zhu.pdf}.

\end{thebibliography}

\paragraph{AI-assisted research and writing disclosure.}
Carptopus is the author responsible for the mathematical claims and the public version of this
manuscript. OpenAI Codex was used as an AI-assisted tool for literature organization, proof
development and stress testing, exact verification, adversarial auditing, and manuscript
preparation.

\end{document}
